\begin{frame}{세 행렬의 ``역할''은 학습이 만든다}
  2단어 예의 세 행렬이 각각 하는 일:
  \small
  \begin{tabular}{@{}llp{8cm}@{}}
    \toprule
    행렬 & 효과 & 역할 \\
    \midrule
    $W_Q$ & 두 번째 차원을 2배로 확대 & ``찾는 방향''을 세우다 \\
    $W_K$ & 모든 차원을 절반으로 축소 & ``색인''의 스케일을 줄이다 \\
    $W_V$ & 두 차원을 서로 섞음 (비대각 원소 $\neq 0$) &
      전달할 \textbf{내용}을 다른 공간에서 구성하다 \\
    \bottomrule
  \end{tabular}
  \vskip0.5em
  이렇게 서로 다른 ``역할''을 하는 변환들은 \textbf{학습을 통해
  만들어진다} — 사람이 ``이 행렬은 대명사 담당''이라고 설계하는
  것이 아니다.
  \vskip0.3em
  {\footnotesize 실습: \texttt{qkv(X, WQ, WK, WV)} 함수를 완성해
  손계산한 $Q, K, V$ 와 일치하는지 확인 — ``같은 입력에서 세 가지
  다른 결과''가 나옴이 이 절의 포인트.}
\end{frame}
